% Datasheet 对比报告模板 —— XeLaTeX 编译两遍
% 占位符 <...> 全部替换；不需要的节删掉，节顺序不要变
\documentclass[11pt,a4paper]{ctexart}
\usepackage[margin=2.0cm]{geometry}
\usepackage{booktabs,longtable,array,multirow}
\usepackage[table]{xcolor}
\usepackage{pdflscape}
\usepackage{tikz}
\usepackage{hyperref}
\hypersetup{colorlinks=true,linkcolor=blue!60!black,urlcolor=blue!60!black}
\definecolor{hdr}{RGB}{31,78,121}
\definecolor{rowa}{RGB}{240,245,250}
\definecolor{barb}{RGB}{47,111,179}
\definecolor{barr}{RGB}{184,58,58}
\definecolor{trk}{RGB}{236,239,243}
\newcommand{\best}[1]{\textcolor{green!45!black}{\textbf{#1}}}
\newcommand{\worst}[1]{\textcolor{red!70!black}{\textbf{#1}}}
\newcommand{\nd}{--}
\newcommand{\hd}[1]{\textcolor{white}{\textbf{#1}}}
\renewcommand{\arraystretch}{1.3}

\title{\textbf{<主题> N 款竞品对比报告}\\[4pt]
\large <型号A> / <型号B> / <型号C>}
\author{基于目录内 N 份原厂规格书整理}
\date{<YYYY年M月D日>}

\begin{document}
\maketitle

\section{概述与结论}
本报告对比 N 款 <器件类别>。除特别注明外，所有参数取 <统一测试条件> 的典型值。
\begin{itemize}
  \item \textbf{<维度1>最优}：<型号>，<数值>；其余 <范围>。
  \item \textbf{<维度2>最优}：<型号>，<数值>。
\end{itemize}
\textbf{面向 <目标应用> 的选型判断。}<一段话：先给结论，再给依据，点出换料风险>。

\section{样本说明}
\begin{longtable}{>{\raggedright\arraybackslash}p{3.8cm}>{\raggedright\arraybackslash}p{3.0cm}>{\raggedright\arraybackslash}p{2.4cm}>{\raggedright\arraybackslash}p{2.4cm}>{\raggedright\arraybackslash}p{3.4cm}}
\toprule
\rowcolor{hdr}\hd{文件} & \hd{型号} & \hd{厂商} & \hd{版本 / 日期} & \hd{备注}\\
\midrule
\endhead
<file\_\allowbreak a.pdf> & <型号A> & <厂商> & <Rev / 日期> & <扫描件 / 预发布 / 分档 等>\\
\rowcolor{rowa}<file\_\allowbreak b.pdf> & <型号B> & <厂商> & <Rev / 日期> & \\
\bottomrule
\end{longtable}

% ---- 参数总表：器件 ≤5 列用竖版；>5 列用下面的横版 ----
\begin{landscape}
\section{参数总表（<统一测试条件>）}
\small
\setlength{\tabcolsep}{4pt}
% 横版版心 25.7cm：首列 2.9cm + 8 列 × 2.4cm + 列距 ≈ 24.6cm
\begin{longtable}{>{\raggedright\arraybackslash}p{2.9cm}*{8}{>{\centering\arraybackslash}p{2.4cm}}}
\toprule
\rowcolor{hdr}\hd{参数} & \hd{<A>} & \hd{<B>} & \hd{<C>} & \hd{<D>} & \hd{<E>} & \hd{<F>} & \hd{<G>} & \hd{<H>}\\
\midrule
\endhead
工作电压 (V) & 2.5--5.5 & \best{2.2--5.5} & & & & & & \\
\rowcolor{rowa}绝对最大电压 (V) & 25 & \worst{6.5} & & & & & & \\
工作电流 典型(最大) mA & 1.8 (2.3) & \best{0.6 (0.9)} & & & & & & \\
\rowcolor{rowa}灵敏度 典型(范围) & & & & & & & & \\
公差 & $\pm18.5\%$ & \best{$\pm7\%$} & & & & & & \\
\rowcolor{rowa}工作温度 ($^\circ$C) & $-40$--105 & \best{$-40$--125} & & & & & & \\
ESD HBM & \nd & 4\,kV & & & & & & \\
\rowcolor{rowa}封装 & SOT-23 & DFN1.6$\times$1.6 & & & & & & \\
\bottomrule
\end{longtable}
\normalsize
绿色为该行最优，红色为需要留意的短板，“--”表示规格书未给出。
\end{landscape}

\section{<派生分析：如极性 / 行程 / 效率 等>}\label{sec:derived}
<说明推算方法，并注明“为本报告推算值，非厂商标称”>

\section{关键参数图}
\begin{center}
\begin{tikzpicture}[x=2.2cm,y=1cm]  % 满刻度 × x ≈ 6--8cm
\node[font=\bfseries\small,anchor=west] at (-1.9,0.5) {<参数名 单位>（<条件>，满刻度 <F>）};
\foreach \name/\v [count=\i] in {<型号A>/3.15, <型号B>/2.7, {<含斜杠的名字>}/2.3}{
  \node[anchor=east,font=\small] at (-0.1,-0.42*\i) {\name};
  \fill[trk]  (0,-0.42*\i-0.13) rectangle (<F>,-0.42*\i+0.13);
  \fill[barb] (0,-0.42*\i-0.13) rectangle (\v,-0.42*\i+0.13);
  \node[anchor=west,font=\scriptsize] at (\v+0.03,-0.42*\i) {\v};
}
\end{tikzpicture}

\small 图 1：<一句话解释图里最值得看的点>。
\end{center}

% ---- 双向条形图（可选：正/负两侧量，如行程、裕量）----
\begin{center}
\begin{tikzpicture}[x=0.005cm,y=1cm]  % 1000 单位 = 5cm
\foreach \name/\dn/\up/\pd/\pu [count=\i] in {<A>/703/185/N/S, <B>/824/110/S/N}{
  \node[anchor=east,font=\small] at (-1120,-0.55*\i) {\name};
  \fill[barr] (-\dn,-0.55*\i-0.15) rectangle (0,-0.55*\i+0.15);
  \fill[barb] (0,-0.55*\i-0.15) rectangle (\up,-0.55*\i+0.15);
  \node[anchor=east,font=\scriptsize] at (-\dn-10,-0.55*\i) {\dn\,(\pd)};
  \node[anchor=west,font=\scriptsize] at (\up+10,-0.55*\i) {\up\,(\pu)};
}
\draw[black!70] (0,-0.2) -- (0,-1.5);
\draw[black!50] (-1000,-1.5) -- (1000,-1.5);
\foreach \t in {-1000,-500,0,500,1000}{\draw[black!50] (\t,-1.5) -- (\t,-1.6); \node[below,font=\scriptsize] at (\t,-1.6) {\t};}
\fill[barr] (-1000,0.35) rectangle (-960,0.55); \node[anchor=west,font=\scriptsize] at (-950,0.45) {<左侧含义>};
\fill[barb] (-50,0.35) rectangle (-10,0.55);    \node[anchor=west,font=\scriptsize] at (0,0.45) {<右侧含义>};
\end{tikzpicture}

\small 图 2：<说明>。
\end{center}

\section{逐款点评}
\begin{description}
  \item[<型号A>（<厂商 / 定位>）] <优点>。<短板>。<规格书特殊事项：版本、分档、未来产品等>。
  \item[<型号B>（<厂商>）] …
\end{description}

\section{规格书疑点与勘误}\label{sec:err}
\begin{enumerate}
  \item \textbf{<型号>}：<问题描述；若新版已修正写“已在 Vx.x 修正”>。
\end{enumerate}

\vspace{1em}
\noindent\small 数据来源：目录 \texttt{<路径>} 内 N 份 PDF；<扫描件说明>。<推算值说明>。
\end{document}
